Un magasin d'ameublement interroge \np{1000} clients ayant réalisé un achat au cours des 12 derniers mois, soit en magasin, soit sur internet (un seul mode d'achat par client).

\medskip

Parmi les \np{1000} clients, 400 ont acheté sur internet.

90\,\% des acheteurs sur internet sont satisfaits et 20\,\% des acheteurs en magasin ne sont pas satisfaits.

\medskip

On choisit au hasard un client et on considère les évènements :
\begin{itemize}[left=5mm]
	\item $I$ : \og le client a acheté sur internet \fg ;
	\item $S$ : \og le client est satisfait \fg ;
\end{itemize}
$\overline{I}$ et $\overline{S}$ sont les évènements contraires respectifs de $I$ et de $S$.

\begin{enumerate}
	\item Compléter le tableau ci-dessous représentant la situation.
	
\begin{center}
	\begin{tblr}{
			colspec = {Q[c,m,1.8in] | Q[c,m,1.8in] | Q[c,m,1.8in] | Q[c,m,1.2in]},
			hlines, vlines,
			row{1} = {font=\bfseries, bg=gray!15},
			column{1} = {font=\bfseries, bg=gray!15}
		}
		& Achat sur internet $\big(I\big)$ & Achat en magasin $\big(\overline{I}\big)$ & Total \\
		Satisfait $\big(S\big)$ & & & \\
		Non satisfait $\big(\overline{S}\big)$ & & & \\
		Total & 400 & & \np{1000} \\
	\end{tblr}
\end{center}	
	
	\item Calculer la probabilité que le client ait acheté en magasin et soit satisfait.
	\item Sachant que le client n'est pas satisfait, calculer la probabilité qu'il ait fait son achat sur internet.
	\item Une stratégie d'amélioration sera mise en place si plus de 10\,\% des clients ne sont pas satisfaits de leur achat.

	Cette stratégie sera-t-elle mise en place ? Justifier.
\end{enumerate}
